\documentclass[11pt]{article}
\usepackage[margin=0.65in]{geometry}
\usepackage{lmodern}
\usepackage{microtype}
\usepackage{enumitem}
\usepackage{tabularx}
\usepackage[table]{xcolor}
\usepackage{titlesec}
\usepackage[hidelinks]{hyperref}
\setlength{\parindent}{0pt}
\setlength{\parskip}{0.22em}
\setlist[itemize]{leftmargin=1.3em,itemsep=0.05em,topsep=0.08em}
\titleformat{\section}{\large\bfseries}{}{0pt}{}
\titlespacing*{\section}{0pt}{0.45em}{0.15em}
\pagestyle{plain}
\newcommand{\blankline}{\par\vspace{0.06em}\noindent\rule{\linewidth}{0.4pt}\par\vspace{0.22em}}
\newcommand{\smallblank}{\rule{0.43\linewidth}{0.4pt}}
\begin{document}
\begin{center}
{\Large\bfseries U1L17SP --- Short Paper 4}\par
{\LARGE\bfseries Who Reasons More Dishonestly with Himself?}\par
{\large Week 17: Literary Self-Deception}
\end{center}
\noindent\textbf{Name:} \smallblank\hfill
\textbf{Date:} \rule{1.35in}{0.4pt}

\begingroup\small
\vspace{0.2em}
\fcolorbox{black}{gray!8}{\begin{minipage}{0.94\linewidth}
\textbf{The question}

Who reasons more dishonestly with himself: Macbeth or Brutus? Compare the private reasoning each man uses before choosing political murder, then defend a fair judgment about the kind and degree of self-deception involved.

You may choose either character, judge them equally self-deceiving in different ways, or argue that the charge fits one of them only partly. Your conclusion must be proportional to the evidence and fair to both characters.
\end{minipage}}

\section*{Text support}
The accompanying \textit{U1L17SP Text Supplement} prints Macbeth's private reasoning in \textit{Macbeth} 1.7 and Brutus's orchard reasoning in \textit{Julius Caesar} 2.1. These are anchors, not limits. You may use any accurate line, action, choice, or consequence from either play; remembered events may be paraphrased accurately.

\section*{The assignment}
Write a \textbf{two- to three-page handwritten comparison} on separate lined paper. Give it the title \textit{Who Reasons More Dishonestly with Himself?} Organize the argument independently. It must:
\begin{itemize}
\item make a comparative claim that defines what \emph{reasoning dishonestly with oneself} means in your argument;
\item reconstruct each character's desire, conclusion, stated reasons, and at least one hidden assumption;
\item apply the logic tools that best expose each case rather than forcing the same diagnosis on both;
\item use at least four brief pieces of evidence, including exact quotations from both supplied passages;
\item give each character fair counterpressure before reaching a qualified comparative judgment.
\end{itemize}
Plot summary is not analysis. Use an event only when you explain how it exposes, supports, or complicates a character's private reasoning.

\section*{The closed-world rule}
You may use only the supplied text supplement, the assigned plays or course-provided primary texts, course handouts, and notes you made during class.

\textbf{Do not use AI, the internet, online summaries, outside criticism, a tutor, a parent, a friend, or anyone else's planning, revising, correcting, sentences, or ideas.}

At the end of your paper, write and sign: \emph{I wrote this paper without outside help.}

\section*{What counts --- 15 points}
\rowcolors{2}{gray!8}{white}
\begin{tabularx}{\linewidth}{>{\bfseries}p{0.25\linewidth}Xr}
Comparative claim & Defines the standard and makes a proportional judgment about both men. & 3\\
Logical autopsy & Reconstructs and diagnoses both private arguments with fitting tools. & 4\\
Textual evidence & Uses accurate, balanced evidence from both supplied passages and the plays. & 3\\
Counterpressure & Treats the strongest complication for each character fairly. & 2\\
Organization & Builds a clear comparison rather than two disconnected summaries. & 2\\
Completion & Is legible, handwritten, and within the assigned length. & 1\\
\end{tabularx}
\endgroup

\newpage
\begin{center}
{\Large\bfseries Minimal Planning Page}\par
{\large U1L17SP --- Make the comparison your own}
\end{center}
\textbf{Do not turn this page into the paper.} Record only the decisions you need before drafting on separate lined paper.

\section*{1. Set a fair standard}
In one sentence, define what counts as reasoning dishonestly with oneself. Your standard must be usable for both Macbeth and Brutus.
\blankline
\blankline

\section*{2. Run two brief autopsies}
For \textbf{Macbeth}, jot: desire; conclusion or intended action; what he admits; hidden bridge; best exposing tool; fair counterpressure.
\blankline
\blankline
\blankline

For \textbf{Brutus}, jot the same six parts. The exposing tool may be different.
\blankline
\blankline
\blankline

\section*{3. Build the comparison}
What is the most important similarity? What difference matters most to your judgment?
\blankline
\blankline

Write the exact quotation from each passage that carries the most weight. Add one broader-play detail for each man, paraphrased accurately if it is not printed in your materials.
\blankline
\blankline
\blankline

\section*{4. State and qualify the judgment}
Write your present comparative claim. Then name what a fair reader could say against it and revise the claim if needed.
\blankline
\blankline
\blankline

\textbf{Possible structures, not requirements:} point-by-point comparison; one standard tested against both men; apparent winner followed by decisive complication. Choose the structure that keeps the comparison active throughout.

\vfill
\textbf{Before submitting:} Underline your comparative claim. Mark evidence from Macbeth with M and evidence from Brutus with B. Sign the authorship statement.
\end{document}
